\begin{statementbox}
\textbf{条件：}设右开口抛物线
\[
 C:\ y^2=2px\qquad(p>0),
\]
其焦点为 $F\left(\frac p2,0\right)$。固定点
\[
 \boxed{A\left(\frac p2,p\right)}\in C.
\]
点 $M,N$ 是抛物线上的两个\textbf{互不相同}的动点，且均不与 $A$ 重合；直线 $AM,AN$ 的斜率存在。

\begin{center}
 \begin{tikzpicture}[font=\small,>=stealth,line join=round]
\draw[gray!45,rounded corners=3pt] (0,0) rectangle (14.2,9.4);
\node at (7.1,8.8) {两条弦斜率相反，第三条弦的方向固定};
\begin{scope}[shift={(2.4,5.0)},scale=.85]
\draw[->,gray] (-.65,0)--(5.5,0) node[above] {$x$};
\draw[->,gray] (0,-4.8)--(0,3.3) node[left] {$y$};
\draw[teal!65!black,thick,samples=120,domain=-4.5:3.1,variable=\t] plot ({\t*\t/4},\t);
\draw[blue!75!black,very thick] (0,0)--(1,2);
\draw[orange!85!black,very thick] (1,2)--(4,-4);
\draw[green!45!black,very thick] (0,0)--(4,-4);
\fill (1,2) circle(2pt) node[above left=10pt] {$A$};
\fill (0,0) circle(2pt) node[above left=9pt] {$M=O$};
\fill (4,-4) circle(2pt) node[above right=12pt] {$N$};
\fill (1,0) circle(1.8pt) node[below right=10pt] {$F$};
\end{scope}
\node[anchor=west,align=left] at (8.7,5.3)
{$y^2=4x,\quad p=2$\\[12pt]$A(1,2),\quad F(1,0)$\\[12pt]
$M(0,0),\quad N(4,-4)$\\[16pt]
$\textcolor{blue!75!black}{k_{AM}=2},\quad\textcolor{orange!85!black}{k_{AN}=-2}$\\[16pt]
$\textcolor{green!45!black}{k_{MN}=-1}$};
\node[text width=13.2cm,align=center] at (7.1,.65) {坐标轴单位长度相同；斜率相反约束的是两条直线，并非两点坐标。};
\end{tikzpicture}

\end{center}

\textbf{核心结论：}若两条连线的斜率互为相反数，则第三条弦的斜率恒定：
\[
 \boxed{k_{AM}=-k_{AN}\quad\Longrightarrow\quad k_{MN}=-1.}
\]
\end{statementbox}
\clearpage
\begin{statementbox}
\textbf{速查公式：}令 $M=(\frac{u^2}{2p},u)$，$N=(\frac{v^2}{2p},v)$，则
\[
 \boxed{k_{AM}=\frac{2p}{p+u}},\qquad
 \boxed{k_{AN}=\frac{2p}{p+v}},\qquad
 \boxed{k_{MN}=\frac{2p}{u+v}}.
\]
以上各式仅在对应的斜率存在、两点不重合时使用。由题设立刻得到
\[
 \boxed{u+v=-2p}\quad\Longrightarrow\quad\boxed{k_{MN}=-1}.
\]
\textbf{适用边界：}这里的“斜率相反”不是“两个点纵坐标相反”；若 $AM$ 或 $AN$ 竖直、两点重合，就不能使用题设的有限斜率关系。
\end{statementbox}


\clearpage
\begin{statementbox}
\textbf{补充结论：从 C097 得到 C098 的一般形式}

设 $A_0=\left(\frac{a^2}{2p},a\right)$ 是抛物线 $y^2=2px$ 上任意一个\textbf{非顶点}定点（$a\ne0$），$M,N$ 是抛物线上另外两个互不相同的点，$k_{A_0M},k_{A_0N}$ 都存在。若
\[
 k_{A_0M}=-k_{A_0N},
\]
则
\[
 \boxed{\,k_{MN}=-\frac12k_{OA_0}=-\frac pa\,}.
\]
原题取 $A_0=A=(p/2,p)$，即 $a=p$，有 $k_{OA}=2$，于是
\[
 \boxed{k_{MN}=-\frac12\cdot2=-1}.
\]
\textbf{注意：}以上推广要求 $A_0\ne O$，因为 $A_0=O$ 时 $k_{OA_0}$ 不存在；$M$ 或 $N$ 可以是顶点 $O$，需用坐标法补充验证。
\end{statementbox}
